% Fichier engendré par tools/generate-tex.py à partir de 09-probabilites-et-statistiques.
% Les démonstrations en français sont transcrites du fichier Lean (skill transcribe-lean-proof).
\documentclass[11pt,a4paper]{article}

% Compile aussi bien avec pdflatex qu'avec xelatex ou tectonic.
\usepackage{iftex}
\ifPDFTeX
  \usepackage[T1]{fontenc}
  \usepackage[utf8]{inputenc}
\else
  \usepackage{fontspec}
\fi
\usepackage[french]{babel}
\usepackage{amsmath,amssymb,amsthm}
\usepackage[margin=2.5cm]{geometry}
\usepackage[hidelinks]{hyperref}

\theoremstyle{definition}
\newtheorem{definition}{Définition}
\theoremstyle{plain}
\newtheorem{theoreme}{Théorème}
\newtheorem{lemme}[theoreme]{Lemme}

% Un énoncé écrit mais non démontré : la preuve Lean est un `sorry`.
\newcommand{\admis}{\par\smallskip\noindent\textit{Admis : l'énoncé est écrit en Lean, sa démonstration reste à faire.}\par}

% Renvoi à la déclaration Lean, une ligne après la démonstration, aligné à droite.
\newcommand{\source}[2]{\par\smallskip\noindent\hfill{\footnotesize\href{#1}{\texttt{#2}}}\par\smallskip}

\title{Probabilités et statistiques}
\date{}

\begin{document}
\maketitle
% Les identifiants Lean en machine à écrire ne se coupent pas : on tolère des
% espacements plus lâches plutôt que des lignes qui débordent.
\sloppy

\section{ProbabilitesEtStatistiques}

\noindent
Lycée \textemdash{} section \og{} Probabilités et statistiques \fg{}, de la seconde à la
terminale S.

\medskip\noindent
C'est le chapitre où l'écart entre le cours et la formalisation est le plus large. Une
probabilité y est une mesure de masse totale $1$ sur un espace mesurable ; un événement,
une partie mesurable ; l'espérance, une intégrale. Rien de tout cela n'existe au lycée, où
les lois continues sont introduites par leur densité et les théorèmes limites simplement
admis.

\medskip\noindent
Les séries statistiques, elles, sont des familles finies $x : \{0, \dots, n-1\} \to
\mathbb{R}$, dont ce fichier pose lui-même la moyenne, la variance et l'écart-type.

\medskip\noindent
\textbf{Cinq énoncés du programme ne sont traités qu'en partie et deux ne le sont pas du
tout.} Ils sont signalés à leur place, et la note finale récapitule pourquoi : ce sont eux
qui mesurent la distance entre le programme et ce que Mathlib sait aujourd'hui.

\subsection{Définitions : séries statistiques}

\noindent
Toutes les définitions du fichier sont posées ici, avant les démonstrations.

\begin{definition}
La \emph{moyenne} d'une série de $n$ valeurs est
\[ \bar x = \frac{1}{n}\sum_{i} x_i . \]
\end{definition}
\source{https://github.com/Commutator-IO/learn-lean/blob/main/courses/02-lycee/09-probabilites-et-statistiques/ProbabilitesEtStatistiques.lean\#L24}{ProbabilitesEtStatistiques.lean\#L24}

\begin{definition}
La \emph{variance} est la moyenne des carrés des écarts à la moyenne,
\[ V(x) = \frac{1}{n}\sum_{i} (x_i - \bar x)^{2} . \]
\end{definition}
\source{https://github.com/Commutator-IO/learn-lean/blob/main/courses/02-lycee/09-probabilites-et-statistiques/ProbabilitesEtStatistiques.lean\#L27}{ProbabilitesEtStatistiques.lean\#L27}

\begin{definition}
L'\emph{écart-type} est la racine carrée de la variance,
$\sigma(x) = \sqrt{V(x)}$.
\end{definition}
\source{https://github.com/Commutator-IO/learn-lean/blob/main/courses/02-lycee/09-probabilites-et-statistiques/ProbabilitesEtStatistiques.lean\#L31}{ProbabilitesEtStatistiques.lean\#L31}

\begin{definition}
La \emph{médiane} d'une série ordonnée de taille impaire $2m + 1$ est sa valeur centrale,
\[ \mathrm{Med}(x) = x_m . \]

\medskip\noindent
Les \emph{quartiles} et l'\emph{écart interquartile} ne sont pas définis dans le fichier :
ils demandent, comme la médiane, de supposer la série ordonnée, et leur définition scolaire
varie d'un manuel à l'autre selon la convention d'arrondi des rangs.
\end{definition}
\source{https://github.com/Commutator-IO/learn-lean/blob/main/courses/02-lycee/09-probabilites-et-statistiques/ProbabilitesEtStatistiques.lean\#L34}{ProbabilitesEtStatistiques.lean\#L34}

\subsection{Moyenne, médiane, écart-type}

\begin{theoreme}
La médiane partage la série ordonnée : si $x$ est croissante, alors
\[ i \le m \implies x_i \le \mathrm{Med}(x) , \qquad
   i \ge m \implies \mathrm{Med}(x) \le x_i . \]
\end{theoreme}

\begin{proof}
La médiane est $x_m$ par définition, et $x$ est croissante : de $i \le m$ on tire
$x_i \le x_m$, et de $m \le i$ on tire $x_m \le x_i$. C'est la croissance appliquée deux
fois, et c'est tout ce que dit la médiane d'une série ordonnée.
\end{proof}
\source{https://github.com/Commutator-IO/learn-lean/blob/main/courses/02-lycee/09-probabilites-et-statistiques/ProbabilitesEtStatistiques.lean\#L40}{ProbabilitesEtStatistiques.lean\#L40}

\begin{theoreme}
Linéarité de la moyenne :
\[ \overline{ax + b} = a\,\bar x + b . \]
\end{theoreme}

\begin{proof}
On développe la somme :
\[ \frac{1}{n}\sum_i (a x_i + b)
   = \frac{1}{n}\Bigl(a\sum_i x_i + nb\Bigr)
   = a\,\frac{\sum_i x_i}{n} + b , \]
la somme des $n$ constantes $b$ valant $nb$. La division par $n$ réclame $n \ne 0$ : une
série vide n'a pas de moyenne, et l'hypothèse figure dans l'énoncé.
\end{proof}
\source{https://github.com/Commutator-IO/learn-lean/blob/main/courses/02-lycee/09-probabilites-et-statistiques/ProbabilitesEtStatistiques.lean\#L46}{ProbabilitesEtStatistiques.lean\#L46}

\begin{theoreme}
Variance d'une série transformée :
\[ V(ax + b) = a^{2}\,V(x) . \]
\end{theoreme}

\begin{proof}
D'après la linéarité de la moyenne, la moyenne de la série transformée est
$a\bar x + b$. Chaque écart devient donc
\[ (a x_i + b) - (a\bar x + b) = a\,(x_i - \bar x) , \]
dont le carré est $a^{2}(x_i - \bar x)^{2}$. Le facteur $a^{2}$ sort de la somme et de la
division par $n$.

\medskip\noindent
La translation $b$ a disparu : déplacer toute la série ne change pas sa dispersion.
\end{proof}
\source{https://github.com/Commutator-IO/learn-lean/blob/main/courses/02-lycee/09-probabilites-et-statistiques/ProbabilitesEtStatistiques.lean\#L54}{ProbabilitesEtStatistiques.lean\#L54}

\begin{theoreme}
Écart-type d'une série transformée :
\[ \sigma(ax + b) = |a|\,\sigma(x) . \]
\end{theoreme}

\begin{proof}
La variance étant multipliée par $a^{2}$, l'écart-type est multiplié par
$\sqrt{a^{2}} = |a|$. Le passage $\sqrt{a^{2}V} = \sqrt{a^{2}}\sqrt{V}$ demande que $V$
soit positive, ce qu'elle est comme moyenne de carrés.

\medskip\noindent
C'est la valeur absolue, et non $a$, qui apparaît : une dilatation de rapport $-2$
disperse autant qu'une dilatation de rapport $2$.
\end{proof}
\source{https://github.com/Commutator-IO/learn-lean/blob/main/courses/02-lycee/09-probabilites-et-statistiques/ProbabilitesEtStatistiques.lean\#L66}{ProbabilitesEtStatistiques.lean\#L66}

\begin{theoreme}
Formule de König--Huygens :
\[ V(x) = \overline{x^{2}} - \bar x^{2} . \]
\end{theoreme}

\begin{proof}
On développe le carré de l'écart :
\[ (x_i - \bar x)^{2} = x_i^{2} - 2\bar x\,x_i + \bar x^{2} , \]
puis on somme sur les $n$ indices :
\[ \sum_i (x_i - \bar x)^{2}
   = \sum_i x_i^{2} - 2\bar x \sum_i x_i + n\,\bar x^{2}
   = \sum_i x_i^{2} - 2n\,\bar x^{2} + n\,\bar x^{2} , \]
en remplaçant $\sum_i x_i$ par $n\bar x$. En divisant par $n$ il reste
$\overline{x^{2}} - \bar x^{2}$.
\end{proof}
\source{https://github.com/Commutator-IO/learn-lean/blob/main/courses/02-lycee/09-probabilites-et-statistiques/ProbabilitesEtStatistiques.lean\#L76}{ProbabilitesEtStatistiques.lean\#L76}

\subsection{Probabilité d'une réunion, d'un complémentaire}

\begin{theoreme}
\[ P(A \cup B) + P(A \cap B) = P(A) + P(B) , \qquad P(A) + P(\bar A) = 1 . \]
\end{theoreme}

\begin{proof}
Les deux égalités sont celles de la bibliothèque : l'additivité d'une mesure sur une
réunion, comptée avec l'intersection, et le fait qu'un événement et son contraire
partagent l'univers, dont la probabilité vaut $1$.

\medskip\noindent
La forme du programme, $P(A \cup B) = P(A) + P(B) - P(A \cap B)$, est écrite ici sans
soustraction. Ce n'est pas une coquetterie : les mesures de Mathlib sont à valeurs dans
$[0, +\infty]$, où la soustraction est tronquée. L'égalité additive dit la même chose sans
faire intervenir cette convention.
\end{proof}
\source{https://github.com/Commutator-IO/learn-lean/blob/main/courses/02-lycee/09-probabilites-et-statistiques/ProbabilitesEtStatistiques.lean\#L96}{ProbabilitesEtStatistiques.lean\#L96}

\subsection{Espérance et variance d'une variable aléatoire}

\begin{theoreme}
Pour une variable aléatoire $X$ de carré intégrable :
\[ E(aX + b) = a\,E(X) + b , \qquad V(aX + b) = a^{2}V(X) , \qquad
   V(X) = E(X^{2}) - E(X)^{2} . \]
\end{theoreme}

\begin{proof}
\emph{L'espérance} est l'intégrale de $X$. La linéarité de l'intégrale donne
$E(aX + b) = a E(X) + E(b)$, et l'intégrale de la constante $b$ vaut $b$ parce que la
mesure est de masse totale $1$ : c'est le seul endroit où sert le fait que $P$ est une
probabilité.

\emph{La variance} : la bibliothèque démontre séparément qu'ajouter une constante ne change
pas la variance, et que multiplier par $a$ la multiplie par $a^{2}$. On applique la
première à la variable $aX$, puis la seconde.

\emph{König--Huygens} est cité à la bibliothèque, sous l'hypothèse que $X$ est de carré
intégrable \textemdash{} sans quoi $E(X^{2})$ n'existe pas et l'égalité n'a pas de sens.
\end{proof}
\source{https://github.com/Commutator-IO/learn-lean/blob/main/courses/02-lycee/09-probabilites-et-statistiques/ProbabilitesEtStatistiques.lean\#L108}{ProbabilitesEtStatistiques.lean\#L108}

\begin{theoreme}
La variance est positive, et l'écart-type $\sigma = \sqrt{V}$ vérifie $\sigma^{2} = V$.
\end{theoreme}

\begin{proof}
La variance est une intégrale de carrés, donc positive : c'est le résultat de la
bibliothèque. Sa racine carrée, élevée au carré, la redonne, précisément parce qu'elle est
positive \textemdash{} la racine carrée de Lean, définie sur tout $\mathbb{R}$, vaut zéro
sur les négatifs, et l'identité $(\sqrt{v})^{2} = v$ y serait fausse.
\end{proof}
\source{https://github.com/Commutator-IO/learn-lean/blob/main/courses/02-lycee/09-probabilites-et-statistiques/ProbabilitesEtStatistiques.lean\#L123}{ProbabilitesEtStatistiques.lean\#L123}

\subsection{Loi binomiale}

\begin{theoreme}
Loi binomiale : pour $n$ épreuves de Bernoulli indépendantes de paramètre $p$,
\[ P(X = k) = \binom{n}{k} p^{k} (1-p)^{n-k} . \]
\end{theoreme}

\begin{proof}
C'est le résultat de la bibliothèque, qui définit la loi binomiale comme la loi du cardinal
d'une partie aléatoire de $\{0, \dots, n-1\}$ où chaque élément est retenu avec probabilité
$p$, indépendamment des autres \textemdash{} c'est-à-dire exactement le schéma de Bernoulli
du programme. La formule des coefficients binomiaux en découle par dénombrement des parties
à $k$ éléments.
\end{proof}
\source{https://github.com/Commutator-IO/learn-lean/blob/main/courses/02-lycee/09-probabilites-et-statistiques/ProbabilitesEtStatistiques.lean\#L132}{ProbabilitesEtStatistiques.lean\#L132}

\begin{theoreme}
Identité combinatoire : pour tous réels $p$ et $q$,
\[ \sum_{k=0}^{n} k \binom{n}{k} p^{k} q^{n-k} = n\,p\,(p + q)^{n-1} . \]
\end{theoreme}

\begin{proof}
Pour $n = 0$ la somme se réduit à son terme $k = 0$, qui est nul, et le membre de droite
l'est aussi.

Soit $n = m + 1$. Le terme $k = 0$ étant nul, on réindexe la somme sur $k = j + 1$ :
\[ \sum_{j=0}^{m} (j+1)\binom{m+1}{j+1} p^{j+1} q^{m-j} . \]
La relation entre coefficients binomiaux
\[ (m+1)\binom{m}{j} = \binom{m+1}{j+1}(j+1) \]
permet de remplacer $(j+1)\binom{m+1}{j+1}$ par $(m+1)\binom{m}{j}$, et l'on factorise
$(m+1)p$ :
\[ (m+1)\,p \sum_{j=0}^{m} \binom{m}{j} p^{j} q^{m-j} . \]
La somme restante est le développement du binôme de Newton, c'est-à-dire $(p+q)^{m}$. Il
reste $(m+1)p\,(p+q)^{m}$, ce qui est le membre de droite puisque $(m+1) - 1 = m$.
\end{proof}
\source{https://github.com/Commutator-IO/learn-lean/blob/main/courses/02-lycee/09-probabilites-et-statistiques/ProbabilitesEtStatistiques.lean\#L138}{ProbabilitesEtStatistiques.lean\#L138}

\begin{theoreme}
Espérance d'une loi binomiale :
\[ E(X) = np . \]
\end{theoreme}

\begin{proof}
L'espérance sous une loi binomiale est une somme finie : la bibliothèque donne
\[ E(f) = \sum_{k=0}^{n} \binom{n}{k} p^{k}(1-p)^{n-k}\, f(k) , \]
que l'on applique à $f(k) = k$. C'est exactement la somme de l'identité précédente avec
$q = 1 - p$, qui vaut donc
\[ n\,p\,(p + 1 - p)^{n-1} = n\,p\,1^{n-1} = np . \]

\medskip\noindent
Le programme \emph{admet} ce résultat. Il est ici démontré \textemdash{} et Mathlib ne le
contenait pas : le fichier des lois binomiales en énonce le \emph{proof wanted}, c'est-à-dire
un résultat souhaité mais non encore établi.

\medskip\noindent
\textbf{Ce qui n'est pas démontré.} La variance $np(1-p)$ ne l'est pas. Elle demande
l'identité analogue pour $\sum k(k-1)\binom{n}{k}p^{k}q^{n-k}$, obtenue par le même
réindexage appliqué deux fois, puis König--Huygens. Mathlib la donne également comme
\emph{proof wanted}.
\end{proof}
\source{https://github.com/Commutator-IO/learn-lean/blob/main/courses/02-lycee/09-probabilites-et-statistiques/ProbabilitesEtStatistiques.lean\#L165}{ProbabilitesEtStatistiques.lean\#L165}

\subsection{Probabilités conditionnelles}

\begin{theoreme}
Probabilité conditionnelle et probabilités composées :
\[ P_A(B) = \frac{P(A \cap B)}{P(A)} , \qquad P(A) \times P_A(B) = P(A \cap B) . \]
\end{theoreme}

\begin{proof}
Les deux égalités sont celles de la bibliothèque, qui définit la mesure conditionnée à $A$
par la restriction à $A$ renormalisée. La première est la définition écrite avec l'inverse
de $P(A)$ plutôt qu'une barre de fraction ; la seconde s'en déduit, à ceci près qu'elle
reste vraie même si $P(A) = 0$, où les deux membres sont nuls.
\end{proof}
\source{https://github.com/Commutator-IO/learn-lean/blob/main/courses/02-lycee/09-probabilites-et-statistiques/ProbabilitesEtStatistiques.lean\#L185}{ProbabilitesEtStatistiques.lean\#L185}

\begin{theoreme}
Formule des probabilités totales, sur la partition $\{A, \bar A\}$ :
\[ P(B) = P(A)\,P_A(B) + P(\bar A)\,P_{\bar A}(B) . \]
\end{theoreme}

\begin{proof}
Par la formule des probabilités composées, le membre de droite vaut
$P(A \cap B) + P(\bar A \cap B)$, c'est-à-dire $P(B \cap A) + P(B \setminus A)$. Or une
mesure est additive sur cette décomposition : $B$ est la réunion disjointe de ce qui est
dans $A$ et de ce qui n'y est pas.

\medskip\noindent
C'est le calcul que fait un arbre pondéré : on descend chaque branche en multipliant les
probabilités rencontrées, puis on additionne les branches qui aboutissent à $B$.
\end{proof}
\source{https://github.com/Commutator-IO/learn-lean/blob/main/courses/02-lycee/09-probabilites-et-statistiques/ProbabilitesEtStatistiques.lean\#L192}{ProbabilitesEtStatistiques.lean\#L192}

\subsection{Indépendance}

\begin{theoreme}
Deux événements sont indépendants si et seulement si
\[ P(A \cap B) = P(A)P(B) , \]
et cette propriété passe au contraire : $P(A \cap \bar B) = P(A)P(\bar B)$.
\end{theoreme}

\begin{proof}
La première équivalence est la caractérisation de l'indépendance dans la bibliothèque.

Pour la seconde, on part de deux égalités. D'une part $B$ et $\bar B$ découpent $A$ :
\[ P(A \cap \bar B) + P(A \cap B) = P(A) , \]
et en remplaçant $P(A \cap B)$ par $P(A)P(B)$ grâce à l'hypothèse,
\[ P(A \cap \bar B) + P(A)P(B) = P(A) . \]
D'autre part, en factorisant,
\[ P(A)P(\bar B) + P(A)P(B) = P(A)\bigl(P(\bar B) + P(B)\bigr) = P(A) . \]
Les deux quantités $P(A \cap \bar B)$ et $P(A)P(\bar B)$ ont donc la même somme avec
$P(A)P(B)$ ; comme cette dernière est finie, on peut la simplifier, et elles sont égales.

\medskip\noindent
La simplification n'est pas automatique dans $[0, +\infty]$ : $\infty + x = \infty + y$ n'y
entraîne pas $x = y$. C'est pourquoi la démonstration vérifie que $P(A)P(B)$ n'est pas
infinie.
\end{proof}
\source{https://github.com/Commutator-IO/learn-lean/blob/main/courses/02-lycee/09-probabilites-et-statistiques/ProbabilitesEtStatistiques.lean\#L206}{ProbabilitesEtStatistiques.lean\#L206}

\subsection{Loi uniforme}

\begin{theoreme}
Loi uniforme sur $[a\,;b]$ avec $a < b$ : la densité constante $\frac{1}{b-a}$ est bien
une densité de probabilité, et l'espérance vaut le milieu de l'intervalle :
\[ \int_a^b \frac{dx}{b-a} = 1 , \qquad
   \int_a^b \frac{x}{b-a}\,dx = \frac{a+b}{2} . \]
\end{theoreme}

\begin{proof}
L'intégrale d'une constante sur $[a\,;b]$ vaut cette constante multipliée par la longueur :
\[ (b - a) \times \frac{1}{b-a} = 1 , \]
licite puisque $b - a \ne 0$.

Pour l'espérance, l'intégrale de la fonction identité vaut $\frac{b^{2}-a^{2}}{2}$, d'où
\[ \frac{b^{2}-a^{2}}{2(b-a)} = \frac{(b-a)(b+a)}{2(b-a)} = \frac{a+b}{2} . \]
\end{proof}
\source{https://github.com/Commutator-IO/learn-lean/blob/main/courses/02-lycee/09-probabilites-et-statistiques/ProbabilitesEtStatistiques.lean\#L225}{ProbabilitesEtStatistiques.lean\#L225}

\subsection{Loi exponentielle}

\begin{theoreme}
Loi exponentielle de paramètre $r > 0$ : sa densité est $r e^{-rt}$ pour $t \ge 0$, sa
fonction de répartition vaut $1 - e^{-rt}$, donc
\[ P(X > t) = e^{-rt} , \]
et l'absence de mémoire s'écrit
\[ P(X > s + t) = P(X > s)\,P(X > t) . \]
\end{theoreme}

\begin{proof}
La densité et la fonction de répartition sont celles de la bibliothèque, la seconde
obtenue en primitivant la première : $\int_0^t r e^{-ru}\,du = 1 - e^{-rt}$. La condition
$t \ge 0$ est nécessaire : en deçà, la densité est nulle et la fonction de répartition
aussi.

Comme $P(X > t) = 1 - P(X \le t) = e^{-rt}$, l'absence de mémoire se ramène à l'équation
fonctionnelle de l'exponentielle :
\[ e^{-r(s+t)} = e^{-rs} \times e^{-rt} , \]
qui est la propriété $e^{u+v} = e^{u}e^{v}$ lue à l'envers. C'est là tout le contenu de la
\og{} loi sans vieillissement \fg{} : la fonction de survie est une exponentielle, donc
multiplicative.

\medskip\noindent
\textbf{Ce qui n'est pas démontré.} L'espérance $1/r$ ne l'est pas : c'est une intégrale
sur $[0, +\infty[$, donc impropre, et son calcul demande une intégration par parties sur un
intervalle non borné que ce fichier n'établit pas.
\end{proof}
\source{https://github.com/Commutator-IO/learn-lean/blob/main/courses/02-lycee/09-probabilites-et-statistiques/ProbabilitesEtStatistiques.lean\#L240}{ProbabilitesEtStatistiques.lean\#L240}

\subsection{Loi normale}

\begin{theoreme}
Loi normale $\mathcal{N}(m, v)$ : son espérance vaut $m$ et sa variance $v$.
\end{theoreme}

\begin{proof}
Les deux résultats sont ceux de la bibliothèque, qui construit la loi normale par sa
densité $\frac{1}{\sqrt{2\pi v}}\,e^{-(x-m)^{2}/(2v)}$ et calcule ses deux premiers moments
par les intégrales gaussiennes. La loi normale centrée réduite en est le cas $m = 0$,
$v = 1$.

\medskip\noindent
\textbf{Ce qui n'est pas démontré.} Les intervalles $1\sigma$, $2\sigma$, $3\sigma$ \textemdash{}
$68\,\%$, $95\,\%$, $99{,}7\,\%$ \textemdash{} ne le sont pas : ces valeurs numériques
demandent un encadrement de la fonction d'erreur, dont Mathlib ne fournit pas les
estimations à cette précision.

\medskip\noindent
Le \textbf{théorème de Moivre--Laplace} n'est pas formalisé non plus. Mathlib possède un
théorème central limite, dont Moivre--Laplace est un cas particulier, mais l'écart entre son
énoncé \textemdash{} convergence en loi d'une somme de variables indépendantes \textemdash{}
et la forme scolaire, qui parle d'intervalles et de valeurs numériques, demande un travail
de traduction à part entière.

\medskip\noindent
L'\textbf{intervalle de fluctuation asymptotique} au seuil $95\,\%$ et l'\textbf{intervalle
de confiance} $[f \pm 1/\sqrt{n}]$ reposent tous deux sur Moivre--Laplace, et tombent avec
lui. Le programme les admet ; ici ils ne sont pas écrits du tout, ce qui est la seule
réponse honnête.
\end{proof}
\source{https://github.com/Commutator-IO/learn-lean/blob/main/courses/02-lycee/09-probabilites-et-statistiques/ProbabilitesEtStatistiques.lean\#L252}{ProbabilitesEtStatistiques.lean\#L252}

\subsection{Énoncés admis}

\noindent
Ce que le programme demande et que ce document ne démontre pas. Les énoncés sont écrits en Lean, pour que le manque se compte ; leur démonstration est admise.

\begin{theoreme}
La variance d'une loi binomiale de paramètres $n$ et $p$ vaut $np(1 - p)$.
\end{theoreme}
\admis
\source{https://github.com/Commutator-IO/learn-lean/blob/main/courses/02-lycee/09-probabilites-et-statistiques/ProbabilitesEtStatistiques.lean\#L276}{ProbabilitesEtStatistiques.lean\#L276}

\begin{theoreme}
L'espérance d'une loi exponentielle de paramètre $r$ vaut $\tfrac1r$.
\end{theoreme}
\admis
\source{https://github.com/Commutator-IO/learn-lean/blob/main/courses/02-lycee/09-probabilites-et-statistiques/ProbabilitesEtStatistiques.lean\#L281}{ProbabilitesEtStatistiques.lean\#L281}

\begin{theoreme}
\emph{Théorème de Moivre--Laplace.} Pour une loi binomiale de paramètres $n$ et $p$, la
probabilité que la variable centrée réduite $\frac{X_n - np}{\sqrt{np(1-p)}}$ soit inférieure
à $x$ tend, quand $n$ grandit, vers la valeur en $x$ de la fonction de répartition de la loi
normale centrée réduite.
\end{theoreme}
\admis
\source{https://github.com/Commutator-IO/learn-lean/blob/main/courses/02-lycee/09-probabilites-et-statistiques/ProbabilitesEtStatistiques.lean\#L287}{ProbabilitesEtStatistiques.lean\#L287}

\begin{theoreme}
Pour une loi normale d'espérance $\mu$ et d'écart-type $\sigma$, les intervalles
$[\mu \pm \sigma]$, $[\mu \pm 2\sigma]$ et $[\mu \pm 3\sigma]$ ont pour probabilités
respectives $0{,}68$, $0{,}95$ et $0{,}997$, à $0{,}005$ près.
\end{theoreme}
\admis
\source{https://github.com/Commutator-IO/learn-lean/blob/main/courses/02-lycee/09-probabilites-et-statistiques/ProbabilitesEtStatistiques.lean\#L294}{ProbabilitesEtStatistiques.lean\#L294}

\begin{theoreme}
\emph{Intervalle de fluctuation asymptotique au seuil de 95 \%.} La probabilité que la
fréquence observée tombe dans $\left[p \pm 1{,}96\sqrt{\tfrac{p(1-p)}{n}}\right]$ tend vers
celle que la loi normale centrée réduite donne à l'intervalle $[-1{,}96\,;\,1{,}96]$.
\end{theoreme}
\admis
\source{https://github.com/Commutator-IO/learn-lean/blob/main/courses/02-lycee/09-probabilites-et-statistiques/ProbabilitesEtStatistiques.lean\#L303}{ProbabilitesEtStatistiques.lean\#L303}

\begin{theoreme}
\emph{Intervalle de confiance.} À partir d'un certain rang, la fréquence observée est à
moins de $\tfrac{1}{\sqrt n}$ de la probabilité cherchée dans au moins $95\,\%$ des cas.
\end{theoreme}
\admis
\source{https://github.com/Commutator-IO/learn-lean/blob/main/courses/02-lycee/09-probabilites-et-statistiques/ProbabilitesEtStatistiques.lean\#L311}{ProbabilitesEtStatistiques.lean\#L311}

\vfill
\noindent\rule{0.35\linewidth}{0.4pt}\par
\noindent{\footnotesize Leonardo de Moura et Sebastian Ullrich,
\emph{The Lean~4 Theorem Prover and Programming Language},
\emph{Automated Deduction — CADE~28}, Springer, 2021, p.~625--635,
\href{https://doi.org/10.1007/978-3-030-79876-5_37}{doi:10.1007/978-3-030-79876-5\_37}.\par}

\end{document}
